
\documentclass{article}
\usepackage{colortbl}
\usepackage{makecell}
\usepackage{multirow}
\usepackage{supertabular}

\begin{document}

\newcounter{utterance}

\centering \large Interaction Transcript for game `cladder', experiment `full\_v1.5\_default', episode 8207 with score{-}Qwen3.5{-}9B{-}sp.
\vspace{24pt}

{ \footnotesize  \setcounter{utterance}{1}
\setlength{\tabcolsep}{0pt}
\begin{supertabular}{c@{$\;$}|p{.15\linewidth}@{}p{.15\linewidth}p{.15\linewidth}p{.15\linewidth}p{.15\linewidth}p{.15\linewidth}}
    \# & \multicolumn{2}{c}{Player} && \multicolumn{2}{c}{Game Master} \\
    \hline

    \theutterance \stepcounter{utterance}  
    & & & \multicolumn{4}{p{0.6\linewidth}}{
        \cellcolor[rgb]{0.9,0.9,0.9}{
            \makecell[{{p{\linewidth}}}]{
                \texttt{\tiny{[P1$\langle$GM]}}
                \texttt{You are an expert in causal inference. The following question is not a typical commonsense query, but rather a meticulously designed question created by a professor specializing in causal inference, intended to assess the students' mastery of the course content.} \\
\texttt{Imagine a self{-}contained, hypothetical world with only the following conditions, and without any unmentioned factors or causal relationships: Tijv has a direct effect on tijw and xevo. Xevo has a direct effect on gyzp. Tijw has a direct effect on gyzp. For those who are not tijv and are not xevo, the probability of gyzp is 29\%. For those who are not tijv and are xevo, the probability of gyzp is 14\%. For those who are tijv and are not xevo, the probability of gyzp is 11\%. For those who are tijv and are xevo, the probability of gyzp is 16\%. The overall probability of tijv is 88\%. Will xevo increase the chance of gyzp?} \\
\texttt{Start your answer with 'Yes' or 'No', followed by additional reasoning or evidence to support your explanation.} \\
            }
        }
    }
     \\ \\

    \theutterance \stepcounter{utterance}  
    & \multicolumn{4}{p{0.6\linewidth}}{
        \cellcolor[rgb]{0.561,0.737,0.561}{
            \makecell[{{p{\linewidth}}}]{
                \texttt{\tiny{[P1$\rangle$GM]}}
                \texttt{No} \\
\\ 
\texttt{**Reasoning:**} \\
\\ 
\texttt{To determine whether \$xevo\$ increases the chance of \$gyzp\$, we must calculate the conditional probability \$P(gyzp | xevo)\$ and compare it to the marginal probability \$P(gyzp)\$ (or equivalently, compare \$P(gyzp | xevo)\$ to \$P(gyzp | \textbackslash{}neg xevo)\$).} \\
\\ 
\texttt{**1. Define the Variables and Given Probabilities:**} \\
\texttt{Let \$T\$ represent \$tijv\$ and \$X\$ represent \$xevo\$. Let \$G\$ represent \$gyzp\$.} \\
\texttt{The problem provides the following conditional probabilities for \$G\$:} \\
\texttt{*   \$P(G | \textbackslash{}neg T, \textbackslash{}neg X) = 0.29\$} \\
\texttt{*   \$P(G | \textbackslash{}neg T, X) = 0.14\$} \\
\texttt{*   \$P(G | T, \textbackslash{}neg X) = 0.11\$} \\
\texttt{*   \$P(G | T, X) = 0.16\$} \\
\\ 
\texttt{The marginal probability of \$T\$ is given:} \\
\texttt{*   \$P(T) = 0.88\$} \\
\texttt{*   Therefore, \$P(\textbackslash{}neg T) = 1 {-} 0.88 = 0.12\$} \\
\\ 
\texttt{**2. Calculate the Marginal Probability of \$G\$, \$P(G)\$:**} \\
\texttt{Using the Law of Total Probability, we sum over all combinations of \$T\$ and \$X} \\
            }
        }
    }
    & & \\ \\

    \theutterance \stepcounter{utterance}  
    & & & \multicolumn{2}{p{0.3\linewidth}}{
        \cellcolor[rgb]{0.9,0.9,0.9}{
            \makecell[{{p{\linewidth}}}]{
                \texttt{\tiny{[GM$|$GM]}}
                \texttt{No} \\
            }
        }
    }
    & & \\ \\

    \theutterance \stepcounter{utterance}  
    & & & \multicolumn{2}{p{0.3\linewidth}}{
        \cellcolor[rgb]{0.9,0.9,0.9}{
            \makecell[{{p{\linewidth}}}]{
                \texttt{\tiny{[GM$|$GM]}}
                \texttt{yes} \\
            }
        }
    }
    & & \\ \\

    \theutterance \stepcounter{utterance}  
    & & & \multicolumn{2}{p{0.3\linewidth}}{
        \cellcolor[rgb]{0.9,0.9,0.9}{
            \makecell[{{p{\linewidth}}}]{
                \texttt{\tiny{[GM$|$GM]}}
                \texttt{game\_result = LOSE} \\
            }
        }
    }
    & & \\ \\

\end{supertabular}
}

\end{document}
